\documentclass[14pt,a4paper]{extarticle}

% ── Encoding & Language ───────────────────────────────────────────────────────
\usepackage[utf8]{inputenc}
\usepackage[T5]{fontenc}
\usepackage{lmodern}
\usepackage[vietnamese]{babel}

% ── Layout ────────────────────────────────────────────────────────────────────
\usepackage[top=3cm, bottom=3cm, left=3.5cm, right=2cm]{geometry}
\usepackage{setspace}
\onehalfspacing

% ── Typography ────────────────────────────────────────────────────────────────
\usepackage{microtype}
\usepackage{parskip}
\setlength{\parskip}{6pt}

% ── Graphics & Color ──────────────────────────────────────────────────────────
\usepackage{xcolor}
\usepackage{graphicx}
\usepackage{tikz}

% ── Tables & Lists ────────────────────────────────────────────────────────────
\usepackage{booktabs}
\usepackage{tabularx}
\usepackage{enumitem}
\setlist[itemize]{noitemsep, topsep=4pt}

% ── Code / Mono ───────────────────────────────────────────────────────────────
\usepackage{listings}
\usepackage{fancyvrb}

% ── Boxes ─────────────────────────────────────────────────────────────────────
\usepackage{tcolorbox}
\tcbuselibrary{skins, breakable}

\newtcolorbox{insightbox}[1][]{
  enhanced, breakable,
  colback=blue!4!white,
  colframe=blue!40!black,
  fonttitle=\bfseries,
  title={#1},
  left=6pt, right=6pt, top=4pt, bottom=4pt,
}

\newtcolorbox{uibox}[1][]{
  enhanced, breakable,
  colback=gray!6!white,
  colframe=gray!50!black,
  fonttitle=\bfseries\ttfamily,
  title={#1},
  left=8pt, right=8pt, top=6pt, bottom=6pt,
}

\newtcolorbox{highlightbox}{
  enhanced,
  colback=orange!8!white,
  colframe=orange!60!black,
  left=6pt, right=6pt, top=4pt, bottom=4pt,
  fontupper=\itshape,
}

% ── Hyperlinks ────────────────────────────────────────────────────────────────
\usepackage[hidelinks]{hyperref}

% ── Section style ─────────────────────────────────────────────────────────────
\usepackage{titlesec}
\titleformat{\section}{\normalfont\large\bfseries}{\thesection}{1em}{}
\titleformat{\subsection}{\normalfont\normalsize\bfseries}{\thesubsection}{1em}{}

% ─────────────────────────────────────────────────────────────────────────────
\begin{document}

\begin{center}
  {\Large\bfseries CHƯƠNG 1: GIỚI THIỆU}
\end{center}
\vspace{1em}

% ─────────────────────────────────────────────────────────────────────────────
\section{Bối cảnh và động lực nghiên cứu}

Trong lịch sử phát triển của trí tuệ nhân tạo, cờ vua luôn là một trong những
thử thách biểu tượng nhất --- từ Deep Blue đánh bại Garry Kasparov năm 1997,
đến AlphaZero tự học từ đầu và vượt qua Stockfish chỉ sau 4 giờ huấn luyện
năm 2017. Tuy nhiên, nghịch lý là dù AI đã chinh phục cờ vua ở tầm thế giới,
hàng triệu người chơi bình thường vẫn không có công cụ thực sự phù hợp để cải
thiện trình độ của chính mình.

Lichess.org --- nền tảng cờ trực tuyến mã nguồn mở lớn nhất thế giới --- xử
lý trung bình hơn \textbf{500.000 ván đấu mỗi ngày}, tạo ra một khối lượng dữ
liệu khổng lồ chứa đựng toàn bộ hành vi, sai lầm và pattern tư duy của người
chơi ở mọi cấp độ. Đặc biệt, mỗi ván đấu được ghi nhận đến cấp độ \emph{từng
nước đi} --- bao gồm trạng thái bàn cờ, đồng hồ còn lại tính đến từng
centi-giây, khai cuộc, và kết quả. Dữ liệu này hoàn toàn công khai --- nhưng
gần như chưa được khai thác có hệ thống để phục vụ mục tiêu huấn luyện cá nhân
hóa.

\begin{highlightbox}
Đề tài này xuất phát từ một quan sát đơn giản: \textbf{khoảng cách giữa dữ
liệu đang tồn tại và insight có thể hành động được là rất lớn} --- và đó chính
xác là bài toán mà kỹ thuật dữ liệu và trí tuệ nhân tạo hiện đại có thể giải
quyết.
\end{highlightbox}

% ─────────────────────────────────────────────────────────────────────────────
\section{Đặt vấn đề}

\subsection{Vấn đề của người chơi nghiệp dư}

Một người chơi cờ ở mức rating 1200--1600 trên Lichess thường trải qua một
vòng lặp không tiến bộ: chơi nhiều, thua nhiều, nhưng không biết mình sai ở
đâu. Họ biết kết quả --- thua --- nhưng không có cơ chế phản hồi có cấu trúc
để hiểu \emph{nguyên nhân sâu xa}:

\begin{itemize}
  \item Lỗi sai xảy ra nhất quán ở giai đoạn nào của ván đấu?
  \item Chất lượng nước đi giảm như thế nào khi đồng hồ cạn dần?
  \item Khai cuộc nào đang thực sự bị hiểu sai về mặt chiến lược?
  \item Pattern tư duy nào đang lặp đi lặp lại xuyên suốt hàng trăm ván?
\end{itemize}

Công cụ phân tích hiện tại --- kể cả công cụ tích hợp của Lichess --- chỉ phân
tích \textbf{từng ván riêng lẻ}. Người chơi biết họ mắc blunder ở nước 23 của
ván hôm nay, nhưng không ai tổng hợp và nói với họ rằng: \textit{``trong 3
tháng qua, 71\% blunder của bạn xảy ra ở middlegame khi đồng hồ xuống dưới 20
giây.''}

\subsection{Khoảng trống kỹ thuật}

Từ góc độ hệ thống, bài toán chia thành ba lớp:

\begin{description}[leftmargin=2cm, style=nextline]
  \item[\textbf{Lớp thu thập.}]
    Lichess cung cấp API streaming theo thời gian thực, nhưng không có hệ
    thống nào thu thập liên tục, xử lý và lưu trữ dữ liệu theo cấu trúc phù
    hợp cho phân tích sâu --- đặc biệt là dữ liệu đồng hồ ở cấp độ từng nước
    đi.

  \item[\textbf{Lớp xử lý.}]
    Dữ liệu thô từ Lichess là chuỗi ký hiệu SAN (Standard Algebraic
    Notation) --- cần được giải mã thành từng trạng thái bàn cờ (FEN), gán
    nhãn giai đoạn ván đấu, và liên kết với metadata người chơi trước khi có
    thể phân tích.

  \item[\textbf{Lớp trí tuệ.}]
    Ngay cả khi có dữ liệu sạch, việc chuyển đổi con số thống kê thành lời
    khuyên cụ thể, có ngữ cảnh, bằng ngôn ngữ tự nhiên --- đòi hỏi sự kết hợp
    giữa truy vấn dữ liệu thực tế và mô hình ngôn ngữ lớn (LLM).
\end{description}

\subsection{Câu hỏi nghiên cứu}

\begin{highlightbox}
\textbf{Làm thế nào để thiết kế và triển khai một hệ thống kỹ thuật dữ liệu
end-to-end --- từ thu thập streaming thời gian thực đến phân tích AI --- có khả
năng biến dữ liệu ván cờ thô thành insight cá nhân hóa, giúp người chơi hiểu
và cải thiện trình độ một cách có căn cứ?}
\end{highlightbox}

% ─────────────────────────────────────────────────────────────────────────────
\section{Mục tiêu đề tài}

\begin{enumerate}[label=\textbf{MT\arabic*.}, leftmargin=2.2cm]

  \item \textbf{Pipeline thu thập dữ liệu thời gian thực.}
    Xây dựng hệ thống streaming kết nối trực tiếp với Lichess API, thu thập
    liên tục sự kiện ván đấu (game start, game end, từng nước đi kèm đồng hồ),
    đưa vào message queue (Kafka) và lưu trữ dạng columnar trên object storage.

  \item \textbf{Kho dữ liệu phân tích (Data Warehouse).}
    Xử lý dữ liệu thô qua Apache Spark, tạo bảng phân tích phẳng ở cấp độ
    từng nước đi --- mỗi bản ghi chứa đầy đủ context: trạng thái bàn cờ
    (FEN), thời gian còn lại, thông tin người chơi, khai cuộc, giai đoạn ván,
    kết quả. Lưu trữ định dạng Apache Iceberg trên Polaris REST catalog.

  \item \textbf{Phân tích hành vi người chơi xuyên nhiều ván.}
    Khai thác dữ liệu clock ở cấp từng nước đi để phát hiện các pattern hành
    vi: điểm thoát lý thuyết khai cuộc, ngưỡng áp lực thời gian, phân phối
    sai lầm theo giai đoạn, và dấu hiệu tilt sau chuỗi ván thua.

  \item \textbf{Phân tích ván đấu từng nước với Stockfish.}
    Tích hợp Stockfish engine để đánh giá thế cờ tại từng nước đi, xây dựng
    biểu đồ eval xuyên suốt ván, và sử dụng LLM để xác định khoảnh khắc then
    chốt quyết định kết quả.

  \item \textbf{Hệ thống bài tập cá nhân hóa.}
    Tự động trích xuất các vị trí bàn cờ từ ván thua của người chơi, tái lập
    đúng điều kiện áp lực thời gian, và tổ chức thành chuỗi bài tập tương tác
    có hướng dẫn từng bước.

  \item \textbf{AI Coach tương tác.}
    Xây dựng tác nhân AI kết hợp LLM (Gemini 2.5 Flash qua Vertex AI) với
    function calling để truy vấn dữ liệu thực tế, trả lời câu hỏi tự nhiên
    của người chơi bằng nhận định có căn cứ từ dữ liệu.

\end{enumerate}

% ─────────────────────────────────────────────────────────────────────────────
\section{Phạm vi và giới hạn}

\textbf{Đề tài tập trung vào:}
\begin{itemize}
  \item Người chơi trên nền tảng Lichess.org (dữ liệu công khai, API mở).
  \item Các loại ván có đồng hồ: Bullet, Blitz, Rapid --- vì dữ liệu clock là
    yếu tố phân tích cốt lõi.
  \item Phân tích hành vi và pattern của người chơi nghiệp dư (rating
    800--2000).
\end{itemize}

\textbf{Đề tài không bao gồm:}
\begin{itemize}
  \item Xây dựng chess engine mới hay huấn luyện model AI chơi cờ.
  \item Phân tích ván đấu của Grandmaster hay giải đấu thế giới.
  \item Chức năng chơi cờ trực tiếp với AI trong thời gian thực.
\end{itemize}

% ─────────────────────────────────────────────────────────────────────────────
\section{Các kịch bản sử dụng và thiết kế tương tác}

\subsection{Game Explorer --- Khám nghiệm một ván thua}

Người chơi nhập Game ID vào ô tìm kiếm. Thay vì tải thẳng bàn cờ, UI trình
bày một màn hình tóm tắt ngắn:

\begin{uibox}[Màn hình tóm tắt ván đấu]
\begin{Verbatim}[fontsize=\small]
  Van #RPJr6MMX  |  Blitz 5+0  |  Sicilian Defense
  Ban (Den, 1487)  vs.  Doi thu (Trang, 1512)  |  Thua

  [!] Phat hien 1 blunder nghiem trong o nuoc 26
      The co lat nguoc hoan toan chi sau 1 nuoc di

  [ Xem toan bo van ]   [ Nhay den diem mau chot ]
\end{Verbatim}
\end{uibox}

Khi người chơi chọn \textit{``Nhảy đến điểm mấu chốt''}, bàn cờ mở ra tại
nước trước blunder --- không tô đỏ nước sai ngay mà hiển thị câu hỏi dẫn dắt:

\begin{uibox}[Hướng dẫn tương tác]
\begin{Verbatim}[fontsize=\small]
  Đây là thế cờ ở nước 25. Bạn đang chơi Đen,
  đồng hồ còn 8 giây.

  Nhìn vào bàn cờ — bạn thấy gì đáng chú ý không?

  [ Gợi ý 1 ]  [ Gợi ý 2 ]  [ Xem đáp án ]
\end{Verbatim}
\end{uibox}

Hệ thống cung cấp ba cấp gợi ý theo thứ tự từ mơ hồ đến cụ thể:

\begin{enumerate}
  \item \textit{``Quét nhanh các quân Trắng --- quân nào đang không có quân
    bảo vệ?''}
  \item \textit{``Mã Đen đang ở e4. Nó có thể với tới ô nào trong 1 nước?''}
  \item Hiển thị mũi tên xanh trực tiếp lên nước đúng.
\end{enumerate}

Sau khi người chơi đã hiểu, UI chuyển sang biểu đồ eval xuyên suốt ván và nút
\textbf{``Phân tích với AI''} --- gọi Gemini tổng hợp toàn bộ dữ liệu để xác
định chuỗi sai lầm dẫn đến blunder quyết định.

\subsection{What If Machine --- Cỗ máy thời gian}

Đây là tính năng nổi bật nhất về mặt trải nghiệm người dùng. Từ vị trí blunder
vừa phân tích, người chơi có thể chọn một nước đi khác --- hệ thống lập tức
hiển thị \textbf{hai dòng thời gian song song}: ván đấu thực tế đã xảy ra và
ván đấu lẽ ra có thể xảy ra nếu chọn nước đúng.

\begin{uibox}[Giao diện hai dòng thời gian song song]
\begin{Verbatim}[fontsize=\small]
  +------------------+------------------+
  |  THUC TE DA XAY  |  NEU DI Nxd2...  |
  |  RA              |                  |
  |  [ban co]        |  [ban co]        |
  |  eval: -2.8      |  eval: +1.2      |
  |  Xac suat thang: |  Xac suat thang: |
  |       12%        |       71%        |
  |                  |                  |
  |  -> Trang chieu  |  -> Den bat Xe,  |
  |     Vua, Den thua|     loi vat chat |
  +------------------+------------------+

  [ Tu choi tiep tu nhanh nay ]
\end{Verbatim}
\end{uibox}

Stockfish tự động chơi tiếp cả hai nhánh song song, cho người dùng thấy trực
quan ``tương lai'' của từng lựa chọn. Người chơi có thể click \textit{``Tự chơi
tiếp từ nhánh này''} để cầm Đen trong nhánh nước đúng --- trải nghiệm cảm giác
đang thắng từ vị trí mà trong thực tế họ đã thua, tạo tác động tâm lý và học
tập rất mạnh.

Chỉ số \textit{xác suất thắng} được tính từ eval của Stockfish theo công thức
chuẩn trong lý thuyết chess engine, không phải ước tính tùy ý.

\subsection{Pressure Drill --- Luyện tập dưới áp lực thực tế}

Dữ liệu \texttt{clock\_remaining} trong hệ thống ghi nhận chính xác số giây
còn lại tại mỗi nước đi. Kịch bản này khai thác trường dữ liệu đó để tái lập
đúng áp lực thời gian của ván thật --- không phải luyện với thời gian vô hạn,
mà luyện đúng điều kiện khắc nghiệt đã xảy ra.

\begin{uibox}[Giao diện bài tập có đồng hồ đếm ngược]
\begin{Verbatim}[fontsize=\small]
  Bai 3 / 7  --  Trich tu van ngay 20/04, nuoc 31
  ---------------------------------------------
  [Ban co]                    [!] 12 giay

                   Den di. Tim nuoc tot nhat.

                   [ Goi y ]   [ Bo qua ]
\end{Verbatim}
\end{uibox}

Khi hết giờ mà chưa đi, hệ thống phản hồi:

\begin{uibox}[Phản hồi hết thời gian]
\begin{Verbatim}[fontsize=\small]
  Het gio -- dung nhu van that.

  Khi bi ap luc, nao ban chon nuoc an toan thay
  vi nuoc tot nhat. Nuoc dung la Nxd2 -- bat Xe
  khong duoc bao ve, loi 2 diem vat chat.

  [ Thu lai ]   [ Bai tiep theo ]
\end{Verbatim}
\end{uibox}

Sau mỗi phiên luyện, hệ thống tổng kết hiệu suất và so sánh trực tiếp với kết
quả trong ván thật, cho người chơi thấy mức độ cải thiện cụ thể.

\subsection{Phân tích hành vi xuyên nhiều ván}

Đây là lớp phân tích sâu nhất, khai thác toàn bộ lịch sử ván đấu đã thu thập.
Từ dữ liệu clock ở cấp từng nước đi, hệ thống có thể rút ra các insight không
thể có nếu chỉ lưu kết quả ván:

\begin{insightbox}[Điểm thoát lý thuyết khai cuộc]
Trong opening, người chơi đi nhanh vì đang theo lý thuyết. Có một nước cụ thể
mà thời gian suy nghĩ đột ngột tăng gấp 3--4 lần --- đó là nước họ ``ra khỏi
sách'' và bắt đầu tự tư duy. Hệ thống xác định nước đó và tạo bài tập huấn
luyện \emph{khả năng định hướng chiến lược} ngay tại điểm chuyển tiếp --- kỹ
năng mà hầu hết người nghiệp dư thiếu nhất.
\end{insightbox}

\begin{insightbox}[Ngưỡng áp lực cá nhân]
Mỗi người chơi có một ngưỡng đồng hồ riêng: dưới bao nhiêu giây thì chất
lượng nước đi bắt đầu sụp. Ngưỡng này không giống nhau giữa các người chơi ---
một số người bắt đầu mắc lỗi ở 30 giây, số khác vẫn ổn định đến 10 giây. Hệ
thống xác định ngưỡng cá nhân và tạo bài tập tập trung đúng vào điểm nguy hiểm
đó.
\end{insightbox}

\begin{insightbox}[Winning position syndrome]
Truy vấn các ván mà eval từng vượt +2.0 nhưng kết quả là thua hoặc hòa. Đây là
những ván mà người chơi \emph{đang thắng nhưng không biết cách kết thúc} ---
hoặc bị cạn thời gian đúng lúc cần tính toán nhiều nhất. Bài tập tương ứng bắt
đầu từ thế đang thắng đó và yêu cầu người chơi tự tìm ra con đường chiến thắng.
\end{insightbox}

\begin{insightbox}[Tilt detection --- phát hiện trạng thái mất bình tĩnh]
So sánh thời gian suy nghĩ per move của ván sau chuỗi thua 3 ván liên tiếp với
ván bình thường. Người chơi bị tilt thường đi nhanh hơn 40--60\% --- dấu hiệu
rõ ràng của quyết định cảm xúc thay vì lý trí. Khi phát hiện pattern này, hệ
thống không tạo bài tập mà đưa ra cảnh báo trực tiếp và gợi ý nghỉ ngơi.
\end{insightbox}

% ─────────────────────────────────────────────────────────────────────────────
\section{Đóng góp của đề tài}

Đề tài có ba đóng góp chính:

\textbf{Về kỹ thuật dữ liệu.} Thiết kế và triển khai một pipeline dữ liệu
end-to-end hoàn chỉnh --- từ streaming real-time (Kafka), object storage
(MinIO), batch processing (Apache Spark + Iceberg/Polaris), analytical serving
layer (StarRocks), đến AI agent (Vertex AI) --- trên môi trường Kubernetes. Đây
là kiến trúc hiện đại tương đương các hệ thống data platform ở quy mô
production, được triển khai thực tế chứ không chỉ thiết kế trên lý thuyết.

\textbf{Về ứng dụng AI.} Xây dựng AI agent kết hợp function calling với truy
vấn dữ liệu thực tế, tạo ra câu trả lời dựa trên bằng chứng thay vì kiến thức
chung của LLM --- một pattern ngày càng quan trọng trong các hệ thống AI doanh
nghiệp. Việc tích hợp Stockfish để tính xác suất thắng và so sánh hai dòng
thời gian ván đấu là ứng dụng trực tiếp của AI phân tích vào trải nghiệm người
dùng.

\textbf{Về giá trị người dùng.} Cung cấp cho người chơi nghiệp dư loại insight
mà trước đây chỉ có thể có được từ huấn luyện viên cá nhân --- đặc biệt là các
phân tích dựa trên dữ liệu clock ở cấp từng nước đi, thứ không thể rút ra từ
việc chỉ xem lại ván đấu thông thường.

% ─────────────────────────────────────────────────────────────────────────────
\section{Cấu trúc luận văn}

Phần còn lại của luận văn được tổ chức như sau:

\begin{description}[leftmargin=2.5cm, style=nextline]
  \item[\textbf{Chương 2.}]
    Tổng quan lý thuyết và công trình liên quan --- kiến trúc Lambda/Kappa,
    Apache Iceberg table format, LLM agent với function calling, và các hệ
    thống phân tích cờ hiện có.

  \item[\textbf{Chương 3.}]
    Kiến trúc tổng thể của hệ thống --- thiết kế từng thành phần và luồng dữ
    liệu từ Lichess API đến giao diện người dùng.

  \item[\textbf{Chương 4.}]
    Cài đặt và triển khai --- chi tiết kỹ thuật của streaming pipeline, batch
    processing, AI agent, và môi trường Kubernetes.

  \item[\textbf{Chương 5.}]
    Kết quả thực nghiệm và đánh giá hệ thống.

  \item[\textbf{Chương 6.}]
    Kết luận và hướng phát triển tiếp theo.
\end{description}

\end{document}
